\documentclass[11pt]{article}

\usepackage{amsmath,amssymb}
\usepackage[margin=1in]{geometry}
\usepackage{booktabs}
\usepackage{array}
\usepackage{longtable}
\usepackage{hyperref}

\title{HAOS-IIP Master Synthesis: Program-Level Compression of the Paper Spine}
\author{Tomislav Rupi\'c \\ Independent Researcher}
\date{March 2026}

\begin{document}

\maketitle

\begin{abstract}
This paper is a master synthesis of the current HAOS-IIP archive. It is not a new phase, not a new numerical run, and not a raw concatenation of the existing PDFs. Its job is narrower and more useful: compress the numbered release spine into one readable program document while preserving the bounded status of every claim. The integration rule is simple. Repeated setup appears once, positive results are retained only at the level already frozen in the source papers, superseded variants are marked as archival, and detailed numerics remain delegated to the original release notes and repository artifacts. Across the integrated archive the program establishes: a kernel-induced operator architecture with a validated scalar low-mode branch and explicit vector-sector limits; a frozen pre-geometric regime that supports morphology and collective organization but not robust binding; a negative boundary on weak reinforcement, registration, and delayed-memory extensions; a clustered DK and combinatorial control line that yields texture-level closure without a universal mechanism; a reproducible deterministic readout instrument and a frozen branch-local cochain-Laplacian hierarchy; a bounded emergence ladder from localized persistent modes through collective transport, temporal ordering, causal closure, and proto-geometric distance surrogates; a minimal continuum sketch; and a callable structural-stability oracle. Concrete frozen anchors include Phase V class consistency $1.000000$, Phase VI nullspace estimate $4$ with spectral radius increasing from $7.862310$ to $7.991324$, proto-particle aggregate overlaps $0.999198$ and $0.999423$ at $n_{\mathrm{side}} \in \{48,60\}$, Phase XVII mean edge reproducibility $0.626262626263$ with causal-depth drift $0.0$, Phase XVIII triangle-violation rate $0.0$, and the public oracle ladder baseline/perturbed/fragmented $\to$ stable/marginal/unstable. The bounded conclusion is programmatic: HAOS-IIP has matured into a reproducible research instrument with a coherent architecture, a frozen authority arc, and a small public tool surface. No continuum ontology, metric emergence, particle identification, or physical correspondence is asserted.
\end{abstract}

\tableofcontents

\section{Why A Master Synthesis Exists}

The repository already contains a large paper spine. That spine is valuable, but it creates a practical problem for external readers: the program is legible only if someone knows which releases are foundational, which are capstones, which are boundary notes, and which are tool-facing compressions.

This paper therefore does not try to replace the archive. It provides a higher-level reading layer with three rules:
\begin{enumerate}
\item each result is stated only at the level already justified in the source paper;
\item repeated setup language is collapsed into one common description;
\item the original numbered papers remain authoritative for mechanics, plots, and detailed ledgers.
\end{enumerate}

The result is a program document rather than a bundle of sequential notes.

\section{Compression Rule}

The official source spine is the numbered release stack in \texttt{papers/pdf\_releases/INDEX.md}. This master document compresses that stack into seven blocks.

\begin{center}
\small
\begin{tabular}{p{0.17\linewidth}p{0.20\linewidth}p{0.50\linewidth}}
\toprule
Block & Primary releases & Integrated role \\
\midrule
Architecture & \texttt{00.1--04.1} & Defines the kernel/substrate/operator-lift frame and tests the earliest scalar, vector, and gauge-like diagnostics. \\
Frozen pre-geometry & \texttt{05.1--18.1} & Maps propagation, morphology, collective organization, and the boundary where organization does not become binding. \\
Weak extension boundary & \texttt{19.1--22.2} & Tests weak reinforcement, registration, and delayed-memory variants and closes that nearby mechanism family. \\
DK and control branches & \texttt{23.1--26.1} plus legacy \texttt{27.1} & Explores clustered DK texture, no-distance combinatorial controls, and reduced smeared-transfer closure without claiming universality. \\
Readout and frozen hierarchy & \texttt{28.1--35.1} & Establishes diagnostic authority, freezes the branch-local operator family, and extends the branch to the localized candidate and protection basin. \\
Mesoscopic emergence ladder & \texttt{36.1--44.1} & Builds the coherent chain from interaction and multi-mode sectors to transport, ordering, causal closure, distance surrogates, and a cautious scaling sketch. \\
Tool surface & \texttt{45.1} & Compresses frozen telemetry into a callable structural-stability oracle. \\
\bottomrule
\end{tabular}
\end{center}

One legacy numbered PDF, \texttt{27.1}, exists in the release folder but is not part of the current official index. In this synthesis it is treated as archival glue between the clustered DK closure and the later Phase V authority stack, not as an independent authority layer.

\section{Architecture Layer: What The System Is}

The architecture papers \texttt{00.1--04.1} establish the baseline interpretive rule of the repository:
\[
\text{same kernel}
\quad + \quad
\text{same weighted substrate}
\quad + \quad
\text{different operator lifts}.
\]

That rule matters because it keeps later claims structurally disciplined. The node branch, edge branch, and later Dirac-type or DK-inspired branches are not allowed to come from unrelated microscopic models.

The integrated architectural conclusions are:
\begin{itemize}
\item the scalar branch admits a clean low-mode organization on structured discrete substrates;
\item the edge branch is numerically active but does not yet justify a low coexact vector sector in the strong sense of an emergent gauge field;
\item gauge-like diagnostics are meaningful as bounded operator diagnostics, not as evidence of a full gauge theory;
\item the architecture is strong enough to support later morphology and transport studies without forcing ontological interpretation.
\end{itemize}

This is the first point where the archive is already ``smart'' rather than merely cumulative. Positive results and negative results are both retained. The scalar low-mode branch is a real success. The strong Maxwell-style reading is not.

\section{Frozen Pre-Geometry: Organization Without Binding}

The next large block, \texttt{05.1--18.1}, maps the behavior of a frozen kernel-induced architecture before any weak dynamical reinforcement is introduced.

\subsection{Transport, Morphology, and Collective Structure}

The first half of the block tracks:
\begin{itemize}
\item wave-packet propagation and constraint stability;
\item morphology classes and their perturbation resilience;
\item transition structure between those classes;
\item collective interaction morphology;
\item coarse envelope organization and resolution-interaction coupling.
\end{itemize}

Taken together, these papers show that the frozen branch can support structured organization. Localized patterns survive long enough to be measured. Morphology classes are not random visual artifacts only. Coarse collective envelopes can be defined and compared across controlled settings.

\subsection{Where The Boundary Appears}

The second half of the block asks whether that organization upgrades into something stronger. The papers on effective binding, relational geometry, relational feedback, phase synchronization, and retarded memory all push on that question.

The capstone compression \texttt{18.1} freezes the answer:
\begin{quote}
Frozen pre-geometry can organize, but it does not bind.
\end{quote}

That sentence remains one of the most important global outputs of the archive. It is a negative result, but not a weak one. It sharply marks the limit of what a frozen interaction-invariant substrate can do without new mechanism classes.

\section{Beyond The Boundary: Weak Extensions That Fail}

The Phase II opening papers \texttt{19.1/19.2--22.2} are easy to misunderstand if read quickly. They are not a collapse of the program. They are a disciplined search over the nearest plausible continuation mechanisms beyond the frozen boundary.

The tested families include:
\begin{itemize}
\item weak unfrozen-kernel reinforcement;
\item mesoscopic substrate reinforcement;
\item registration-gated interaction eligibility;
\item delayed nonlocal kernel memory.
\end{itemize}

The integrated conclusion is again structural rather than anecdotal:
\begin{quote}
within this architecture, nearby weak law variants do not generate persistence.
\end{quote}

That matters because it prevents the program from drifting into endless near-neighbor mechanism tweaking. The archive uses these papers to justify a qualitative change of search direction rather than a local parameter chase.

\section{DK And Control Branches: Texture, Austerity, Closure}

The next block, \texttt{23.1--26.1}, opens two important side channels.

First, the DK sector papers study collision texture, grade exchange, and braid-like protection behavior. The strongest positive statement is not that a universal braid mechanism has emerged. The stronger and more honest statement is narrower: under clustered seeds the sector admits a real texture-level organization and a valid effective reduction.

Second, the no-distance combinatorial papers create a deliberate austerity branch. These papers are important methodologically because they strip away metric assumptions and ask what survives under pure combinatorial control structure, incidence noise, motif injection, armed-state irreversibility, and bridge tests.

The clustered DK closure \texttt{26.1} then provides the clean compression:
\begin{itemize}
\item the braid-like signature is not universal on the frozen branch;
\item a valid two-state reduced model does exist in the clustered sector;
\item the stable connected region lies in a smeared-dominant regime;
\item the phase corridor, rather than grade transfer alone, is the primary control variable.
\end{itemize}

This is a classic HAOS-IIP move: keep the reduced description, discard the inflated interpretation.

\section{Readout Authority And The Frozen Branch Hierarchy}

The papers \texttt{28.1--35.1} are where the archive becomes more than a research sequence. They create a reliable authority path.

\subsection{Phase V Readout Authority}

Paper \texttt{28.1} turns the earlier work into a deterministic readout instrument. Its strongest bounded signals are easy to state:
\begin{itemize}
\item class consistency remains fixed at $1.000000$ across a five-schedule reproducibility sweep;
\item the mean-score range is only $0.006224$;
\item stable, degraded, and null controls remain clearly separated.
\end{itemize}

The key conceptual upgrade is that Phase V is not another exploratory theory block. It is a reproducibility instrument.

\subsection{Frozen Operator Family And Spectral Admissibility}

Paper \texttt{29.1} then freezes one branch-valid operator family for later work:
\[
\Delta_h
\quad \text{on} \quad
n_{\mathrm{side}} \in \{12,24,36,48\}.
\]
Every refinement level passes the admissibility checks. The Hermiticity defect is numerically zero, the sampled smallest eigenvalue is nonnegative up to solver precision, the nullspace estimate remains fixed at $4$, and the sampled spectral radius grows monotonically from $7.862310$ to $7.991324$.

This is another crucial transition. Once the operator family is frozen cleanly, later feasibility statements can be compared without moving the substrate.

\subsection{From Spectral Feasibility To The Localized Candidate}

Papers \texttt{30.1--35.1} build upward from that fixed hierarchy:
\begin{itemize}
\item spectral feasibility on the frozen hierarchy;
\item a short-time trace contract;
\item coefficient stabilization and invariant tracking;
\item a cautious continuum-bridge feasibility statement;
\item an integrated proto-particle test;
\item localized-mode protection and failure mechanisms.
\end{itemize}

The integrated proto-particle paper \texttt{34.1} is one of the strongest bounded positive results in the repository. On the added level $n_{\mathrm{side}}=60$, the frozen short-time bridge law stays tight with maximum branch trace prediction error $3.19 \times 10^{-4}$. Two excitation families pass the localization gate at that level, but only one remains aggregate-persistent on both repeat levels. Its overlaps are $0.999198$ and $0.999423$, with recovery-style scores $0.811520$ and $0.887477$, while the matched altered-connectivity control remains diffusive.

Again the claim is bounded. This is proto-particle feasibility, not particle ontology.

\section{Mesoscopic Emergence Ladder}

The papers \texttt{36.1--44.1} form the cleanest integrated narrative in the whole archive. Their logic is not speculative. It is cumulative:
\[
\begin{aligned}
\text{interaction stability}
&\to
\text{multi-mode sectors}
\to
\text{collective transport}
\to
\text{ordering}
\\
&\to
\text{causal closure}
\to
\text{distance surrogate}
\to
\text{cautious scaling sketch}.
\end{aligned}
\]

\subsection{Two-Mode, Multi-Mode, And Collective Transport}

Papers \texttt{36.1--38.1} show that the localized candidate can survive richer composition rules than the single-excitation setting.

By \texttt{38.1}, the dilute sector behaves like a weak mesoscopic medium under bounded probes. The branch exhibits:
\begin{itemize}
\item fixed-layout transport drift $4.96377 \times 10^{-7}$;
\item collective-to-microscopic relaxation ratios between $11.666666666667$ and $35.0$;
\item low-wavenumber fluctuation drift $0.010741131495$;
\item mean identity $0.962058695177$ and mean survival $0.974603174603$ under the maximum weak drive.
\end{itemize}

The altered-connectivity control fails the same stack at least in fluctuation scaling and weak-bias response.

\subsection{Propagation, Ordering, Causal Closure, And Distance}

Paper \texttt{39.1} shows that the frozen collective sector supports bounded propagation structure under deterministic disturbances. On the primary density-pulse probe the branch exhibits effective-speed drift $0.016755292342$, seed spread $0.028668837091$, influence-order mean $0.848677248677$, and low-$k$ transport drift $0.013983435184$, while the deterministic altered-connectivity control fails the same boundedness tests.

Papers \texttt{40.1--42.1} then sharpen the ladder:
\begin{itemize}
\item stable temporal ordering from collective observables;
\item compact directed influence graphs that close into an acyclic causal slice;
\item a depth-based proto-geometric distance surrogate compatible with the same authority slice.
\end{itemize}

The synthesis paper \texttt{43.1} gives the cleanest integrated chain on the frozen authority slice:
\begin{itemize}
\item primary ordering score $0.965986394558$ on the branch;
\item mean edge reproducibility $0.626262626263$;
\item causal-depth drift $0.0$;
\item mean order compatibility $0.96409388653$;
\item triangle-violation rate $0.0$.
\end{itemize}

The matched control broadens materially and fails the same shell-ordering and refinement-scaling gates.

\subsection{The Minimal Continuum Sketch}

Paper \texttt{44.1} is not a new phase. It is a post-freeze compression test over a fixed refinement slice $n_{\mathrm{side}} \in \{48,60,72,84\}$. It evaluates only already frozen observables and asks whether they can be placed on one common refinement axis without losing coherence.

The conclusion is again bounded and useful:
\begin{itemize}
\item several observables show smooth scaling trends under low-complexity fits;
\item the remaining observables stay bounded even where the fits are weaker;
\item the branch stays more regular than the altered-connectivity control on the same axis.
\end{itemize}

This is enough for a proto-continuum feasibility sketch. It is not enough for continuum ontology.

\section{Tool Surface: The First Callable Capability}

Paper \texttt{45.1} matters because it changes the status of the repository. The archive is no longer only a research sequence. It exposes a callable public surface:
\begin{verbatim}
python3 -m haos_iip.demo stability
\end{verbatim}

The tool reuses frozen telemetry for persistence, ordering, acyclicity, causal depth, and order compatibility, then presents the result as a small structural-stability oracle. The default deterministic ladder is:
\begin{center}
\begin{tabular}{lccccc}
\toprule
scenario & retention & consistency & deformation & integrity & class \\
\midrule
baseline   & 1.000 & 1.000 & 0.000 & 1.000 & stable \\
perturbed  & 0.750 & 0.833 & 0.062 & 0.625 & marginal \\
fragmented & 0.250 & 0.656 & 0.172 & 0.375 & unstable \\
\bottomrule
\end{tabular}
\end{center}

The scan mode over a tiny deterministic grid already shows that the surface is not purely static. On the tested baseline grid it yields $3$ stable cells, $0$ marginal cells, and $6$ unstable cells.

This does not promote the program into ontology. It does something more practical: it turns the frozen telemetry layer into a reusable diagnostic capability.

\section{What The Archive Positively Supports}

Taken as one program document, the archive now supports the following bounded positive statements:
\begin{itemize}
\item a coherent operator architecture exists and can be tested on discrete weighted substrates;
\item the frozen architecture supports morphology, transport structure, and collective organization;
\item the program has mapped a real negative boundary: organization does not automatically become binding;
\item weak nearby mechanism families have been explored and broadly closed as persistence-generating candidates;
\item clustered DK and combinatorial control branches support useful reduced descriptions without justifying universal mechanism claims;
\item a deterministic readout instrument and a frozen operator hierarchy provide a reproducible authority path;
\item a bounded emergence ladder from persistent localized modes to temporal ordering, causal closure, and proto-distance surrogates exists on the frozen branch;
\item the public-facing reproduction spine and telemetry oracle make the repository auditable and callable.
\end{itemize}

\section{What The Archive Does Not Support}

The master synthesis is only as good as its non-claims. The integrated archive does \emph{not} justify:
\begin{itemize}
\item continuum limits as established facts;
\item manifold reconstruction;
\item metric emergence in the strong geometric sense;
\item gauge-field emergence in the strong physical sense;
\item particle identification in the physical sense;
\item relativistic light-cone claims;
\item field equations, curvature laws, or cosmological claims;
\item physical ontology.
\end{itemize}

That discipline is not a weakness. It is the main reason the archive remains readable as a research instrument rather than a belief system.

\section{Conclusion}

Yes, the papers can be combined into one master document, but only if the combination is compressive rather than additive. The smart move is not to paste every PDF into one giant stack. The smart move is to preserve the structure of the program:
\begin{itemize}
\item architecture;
\item boundary mapping;
\item failed nearby extensions;
\item control and DK side branches;
\item diagnostic authority;
\item frozen hierarchy;
\item mesoscopic emergence ladder;
\item public tool surface.
\end{itemize}

Seen that way, the archive already reads less like a pile of experiments and more like a scientific instrument with a memory. The individual papers remain necessary. But the program itself can now be seen in one view.

\appendix

\section{Release Provenance Map}

\small
\begin{longtable}{p{0.10\linewidth}p{0.33\linewidth}p{0.47\linewidth}}
\toprule
Release & Short title & Integrated role \\
\midrule
\endfirsthead
\toprule
Release & Short title & Integrated role \\
\midrule
\endhead
\bottomrule
\endfoot
\texttt{00.1} & Operator architecture & Establishes the common kernel/substrate/operator-lift frame and the scalar/vector interpretive baseline. \\
\texttt{01.1} & Kernel-induced transverse band & Shows the emergence of a transverse spectral band on discrete substrates within the kernel-induced architecture. \\
\texttt{02.1} & Continuum reconstruction and scalar-transverse coupling & Tests how the kernel-induced operator approaches continuum-style behavior without claiming continuum ontology. \\
\texttt{03.1} & Dirac-Kaehler factorization & Opens the Dirac-type factorization layer as a structural extension of the operator architecture. \\
\texttt{04.1} & Gauge-like diagnostics & Tests bounded gauge-like observables while explicitly stopping short of a physical gauge claim. \\
\texttt{05.1} & Wave-packet propagation & Establishes early propagation and constraint-stability behavior on the frozen architecture. \\
\texttt{06.1} & Morphology survey & Maps the basic pre-geometric morphology classes supported by the frozen branch. \\
\texttt{07.1} & Perturbation resilience & Tests whether those morphology classes survive bounded perturbations. \\
\texttt{08.1} & Transition structure & Measures the transition logic between morphology classes. \\
\texttt{09.1} & Unified pre-geometric atlas & Compresses the morphology, resilience, and transition stack into one atlas object. \\
\texttt{10.1} & Collective interaction morphology & Begins the collective-organization line on top of the frozen morphology stack. \\
\texttt{11.1} & Coarse-grained envelope structure & Introduces coarse collective envelopes and their reproducible descriptors. \\
\texttt{12.1} & Resolution-interaction coupling & Studies how collective structure changes with resolution and interaction coupling. \\
\texttt{13.1} & Effective binding tests & Searches for weak binding-like behavior and fails to promote the frozen structure into robust binding. \\
\texttt{14.1} & Relational geometry probe & Tests geometry-like relational structure without claiming genuine geometry. \\
\texttt{15.1} & Relational feedback & Examines feedback organization in the same frozen regime. \\
\texttt{16.1} & Phase synchronization probe & Searches for synchronization structure and marks its limitations. \\
\texttt{17.1} & Retarded-memory probe & Tests minimal memory or retardation effects without finding a robust persistence channel. \\
\texttt{18.1} & Frozen pre-geometry limit & Capstone compression of the frozen limit: organization survives, binding does not. \\
\texttt{19.1} & Weak reinforcement opener (superseded) & First release of the unfrozen-kernel and substrate-reinforcement search before \texttt{19.2}. \\
\texttt{19.2} & Weak reinforcement opener & Official Phase II opening paper for weak reinforcement and back-reaction probes. \\
\texttt{20.1} & Mutual registration & Tests registration-gated interaction eligibility as a potential persistence mechanism. \\
\texttt{21.1} & Delayed nonlocal memory & Tests delayed nonlocal memory as a continuation mechanism. \\
\texttt{22.1} & Phase II boundary note & Early compression of the weak-reinforcement, registration, and memory boundary. \\
\texttt{22.2} & Persistence-level null results & Freezes the negative Phase II boundary result at the persistence level. \\
\texttt{23.1} & DK collision and grade exchange & Opens the clustered DK sector and its texture-level interaction baseline. \\
\texttt{23.2} & DK braid topology and protection & Studies braid-like protection structure on the kernel-induced cochain architecture. \\
\texttt{24.1} & Combinatorial austerity branch & Introduces the no-distance edge-graph control branch and its robustness probes. \\
\texttt{25.1} & Extended austerity line & Extends the no-distance line through harmonic-address selection, armed-state irreversibility, and bridge tests. \\
\texttt{25.2} & Austerity closure note & Short closure note for the no-distance combinatorial line. \\
\texttt{25.3} & Full austerity closure & Full closure document for the no-distance combinatorial control branch. \\
\texttt{26.1} & Effective smeared-transfer closure & Compresses the clustered DK sector into a valid reduced closure without universalizing the mechanism. \\
\texttt{27.1} & Legacy Phase IV numbered PDF & Archival transitional object present in the release folder but not part of the current official indexed compression map. \\
\texttt{28.1} & Deterministic readout authority & Freezes the Phase V recovery and control-discrimination instrument. \\
\texttt{29.1} & Frozen operator family & Freezes the branch-local cochain-Laplacian hierarchy and its spectral admissibility. \\
\texttt{30.1} & Spectral feasibility & Establishes bounded spectral-feasibility statements on the frozen hierarchy. \\
\texttt{31.1} & Short-time trace feasibility & Freezes the short-time trace contract that later bridge claims reuse. \\
\texttt{32.1} & Coefficient stabilization & Establishes invariant-tracking and coefficient-stabilization structure. \\
\texttt{33.1} & Cautious continuum bridge & Provides the first bounded continuum-bridge feasibility statement. \\
\texttt{34.1} & Integrated proto-particle feasibility & Establishes the localized candidate under the full frozen authority stack. \\
\texttt{35.1} & Protection and failure mechanisms & Maps the localized candidate's protection basin and breakdown channels. \\
\texttt{36.1} & Two-mode identity preservation & Tests the candidate under two-mode interaction and identity-preservation constraints. \\
\texttt{37.1} & Multi-mode sector formation & Extends the persistent candidate into a multi-mode statistical sector. \\
\texttt{38.1} & Collective mesoscopic transport & Establishes collective dynamics and mesoscopic transport feasibility. \\
\texttt{39.1} & Propagation and causal-scale feasibility & Builds bounded disturbance-spread and causal-scale structure on top of the collective sector. \\
\texttt{40.1} & Temporal ordering & Establishes stable event-ordering structure from frozen collective observables. \\
\texttt{41.1} & Causal closure & Reconstructs a compact acyclic influence graph on the frozen authority slice. \\
\texttt{42.1} & Distance surrogate & Establishes a depth-based proto-geometric distance-surrogate layer. \\
\texttt{43.1} & Ordering, causal closure, and distance synthesis & Compresses Phases XVI--XVIII into one bounded emergent-kinematics stack. \\
\texttt{44.1} & Minimal scaling sketch & Checks whether already frozen observables can be arranged on one common refinement axis coherently. \\
\texttt{45.1} & Structural-stability oracle & Compresses frozen telemetry into a callable public diagnostic tool. \\
\end{longtable}

\normalsize

\end{document}
